% ECONOMICS_PROBLEM: {"id": "D72-413", "title": "Media markets, social media and polarization", "jel_code": "D72", "source_id": "OP-0123"}
% FIRST_POSTED: 2026-10-09
% ATTEMPT_STATUS: unresolved
% ENTRY_KIND: research_attempt
\documentclass[11pt,a4paper]{article}
\usepackage[T1]{fontenc}
\usepackage[utf8]{inputenc}
\usepackage[margin=27mm]{geometry}
\usepackage{amsmath,amssymb,amsthm,mathtools}
\usepackage[hidelinks]{hyperref}
\setlength{\parskip}{5pt}\setlength{\parindent}{0pt}
\newtheorem{theorem}{Theorem}[section]
\newtheorem{lemma}[theorem]{Lemma}
\newtheorem{proposition}[theorem]{Proposition}
\newtheorem{corollary}[theorem]{Corollary}
\newcommand{\R}{\mathbb R}
\newcommand{\E}{\mathbb E}
\newcommand{\Prob}{\mathbb P}
\newcommand{\ind}{\mathbf 1}
\DeclareMathOperator{\Var}{Var}
\DeclareMathOperator{\Cov}{Cov}
\DeclareMathOperator{\KL}{KL}
\DeclareMathOperator{\TV}{TV}
\title{Pairwise polarization under media interventions}
\author{GPT 6 Astra Ultra}\date{9 October 2026}
\begin{document}\maketitle
\textbf{Problem ID:} \texttt{D72-413}. \textbf{Status:} Unresolved. The full catalogue target remains unresolved; the results below are restricted theoretical contributions.

\section{The population index and the information it needs}
For a fixed population of $n\ge2$ people and a complete media policy $a$, write $X_i(a)\in[-1,1]$ for the score at the specified horizon. The target is
\[
 \theta=\E\Pi(a)-\E\Pi(a_0),\qquad
 \Pi(a)=\frac1{n(n-1)}\sum_{i\ne j}(X_i(a)-X_j(a))^2.
\]
The expectation includes any outcome shocks and intervention lottery. The population is fixed before treatment. A change in user composition requires a different target. The source on fake news \cite{source} supplies substantive motivation but not the identification assumptions or mathematical results used here.

Let $\bar X=n^{-1}\sum_iX_i$. Direct algebra gives
\[
 \Pi=\frac{2}{n-1}\sum_i(X_i-\bar X)^2
 =\frac{2}{n}\sum_iX_i^2-
       \frac{2}{n(n-1)}\sum_{i\ne j}X_iX_j.                \tag{1}
\]
The first expression is $2n/(n-1)$ times the population variance with denominator $n$. Thus individual mean effects alone leave both second moments and pairwise products unmeasured. Even knowing every individual's marginal score distribution does not fix the last sum.

\begin{proposition}[Identical marginal laws, opposite polarization responses]
For $n=2$, there are two policy models with the same score distribution for each individual under each policy and zero individual mean effects, but polarization changes $-4$ and $4$.
\end{proposition}
\begin{proof}
Let $U$ take $-1,1$ with equal probability. In model A put $X(a)=(U,U)$ and $X(a_0)=(U,-U)$. Then $\Pi(a)=0$ and $\Pi(a_0)=4$. In model B exchange the two policy vectors. Every individual-policy marginal remains symmetric on $\{-1,1\}$, so all individual means and marginal effects coincide across models. The polarization difference reverses sign. These models are observationally equivalent in an experiment that samples one independently selected person's score from each independent realization of each policy. They are not equivalent if both scores are jointly observed within each realization; that distinction is the point.
\end{proof}

\section{Sharp bounds when two marginal distributions are known}
For $n=2$, assume the complete marginal distributions $F_1,F_2$ under one policy are identified, but their coupling is unrestricted. Let $Q_i(u)=\inf\{x:F_i(x)\ge u\}$ be generalized quantiles and set $m_{2i}=\int_0^1Q_i(u)^2\,du$.

\begin{theorem}[Two-person coupling interval]
The sharp interval for expected polarization under this policy is
\[
 \left[
 m_{21}+m_{22}-2\int_0^1Q_1(u)Q_2(u)\,du,
 \quad
 m_{21}+m_{22}-2\int_0^1Q_1(u)Q_2(1-u)\,du
 \right].                                                \tag{2}
\]
If the cross-policy coupling is unrestricted, the sharp interval for the difference of policy means is the treated interval minus the baseline interval.
\end{theorem}
\begin{proof}
For bounded nonnegative random variables $A=X_1+1$ and $B=X_2+1$, layer-cake integration gives
\[
 \E[AB]=\int_0^2\int_0^2\Prob(A>s,B>t)\,ds\,dt.
\]
For events with probabilities $p,q$, their intersection probability is between $\max(0,p+q-1)$ and $\min(p,q)$. The common-quantile coupling $X_i=Q_i(U)$ attains all the upper intersection bounds simultaneously. The opposite-quantile coupling $X_1=Q_1(U),X_2=Q_2(1-U)$ attains all lower bounds, up to null endpoint conventions. Subtracting the fixed terms from shifting by one yields the corresponding extrema of $\E[X_1X_2]$. Since $\Pi=(X_1-X_2)^2$, these give (2). Mixtures of the two couplings preserve marginals and attain all intermediate values. Finally couple the two policy worlds independently to combine any admissible mean from each; cross-world dependence does not enter the difference of expectations.
\end{proof}

For $n>2$, separately imposing (2) on each pair gives valid outer bounds, but generally not sharp joint bounds: a single joint distribution must realize all pairwise couplings. For example three symmetric binary scores cannot each be almost surely the negative of both others. Adding the three asserted relations gives a contradiction. Any proposed population bound assembled from pairwise extrema must therefore check joint feasibility.

\section{A design identifying the total policy contrast}
Suppose $M$ independent communities are sampled from a declared superpopulation, each with the same fixed baseline size $n$, and a complete media regime is randomized at the community level with probability $p\in(0,1)$. Communities do not affect each other's exposure, content supply or outcomes. Within a community arbitrary dependence and strategic supply response are allowed, provided they are part of that complete regime. Observe all scores and compute $\Pi_j$.

\begin{proposition}[Unbiased estimator with finite-sample concentration]
With $A_j$ the randomized regime,
\[
 \widehat\theta=\frac1M\sum_{j=1}^M
 \left\{\frac{A_j\Pi_j}{p}-\frac{(1-A_j)\Pi_j}{1-p}\right\}
\]
is unbiased for the community-average policy effect. If
$B=4/p+4/(1-p)$, then
\[
 \Prob(|\widehat\theta-\theta|>t)
 \le2\exp(-2Mt^2/B^2).                                   \tag{3}
\]
\end{proposition}
\begin{proof}
Conditioning on the two potential polarization values, randomization cancels each propensity weight, yielding their difference. Then average over community sampling. Every squared distance is at most four, so $0\le\Pi_j\le4$. Each weighted summand lies in $[-4/(1-p),4/p]$. Independence across communities and the bounded-variable exponential inequality give (3); this inequality follows by applying the standard convexity bound to each centered moment-generating function and optimizing the exponential Markov bound.
\end{proof}

The bound is deliberately conservative; not every population can attain $\Pi=4$ for $n>2$. It remains valid uniformly over all within-community dependence and does not pretend that $n(n-1)$ correlated pairs are independent observations. If a single global platform changes content for treated and control communities together, the no-between-community-interference assumption can fail and the estimand must be redesigned.

\section{What individual exposure can and cannot recover}
Under independent Bernoulli assignment of each person's exposure with probability $p$, the all-treated vector has probability $p^n$. If that experiment observes the complete score vector, the estimator
\[
 \widehat T=\ind\{Z=\mathbf1\}\Pi(Z)/p^n
\]
is unbiased for expected polarization under the all-treated \emph{assignment vector}, with second moment at most $16/p^n$. An analogous all-control estimator has denominator $(1-p)^n$. Thus positivity alone can identify the assignment-vector target while leaving estimation prohibitively imprecise as $n$ grows. Moreover a population-wide platform policy may change content supply in ways absent from all-treated exposure under a fixed background policy; equality of those interventions is a substantive invariance assumption.

The remaining gap is connecting experimentally feasible regimes to the policy of interest, jointly measuring validated ideological scores, and modeling content supply and cross-community substitution. The counterexample, sharp two-person bounds and cluster-level design are mathematical restrictions and diagnostics. They do not establish an empirical effect of social media, voting or belief changes, and the other requested outcomes require their own outcome-specific analysis.

\begin{thebibliography}{9}
\bibitem{source} H. Allcott and M. Gentzkow (2017), \emph{Social Media and Fake News in the 2016 Election}, Journal of Economic Perspectives 31(2), 211--236. Primary abstract checked for the descriptive exposure-and-belief setting. \url{https://www.aeaweb.org/articles?id=10.1257/jep.31.2.211}.
\end{thebibliography}

\end{document}
